\ifdefined\gkver\else\def\gkver{student}\fi
\documentclass[\gkver]{gaokaozhenti}
\begin{document}

\chapter{2013年广东理综}
%% paperType: 理综物理部分

\begin{enumerate}
\item
%% number: 13
%% paperName: 2013年广东理综第 13 题 4 分
%% typeId: 1
%% type: 单选题
%% chapter: 直线运动
%% point: 匀变速直线运动
%% method: 无
%% score: 4
%% degree: 820
%% duplicateId: 0
%% body:
某航母跑道长为 $200\Um$，飞机在航母上滑行的最大加速度为 $6\Umsq$，起飞需要的最低速度为 $50\Ums$。那么，飞机在滑行前，需要借助弹射系统获得的最小初速度为\xzanswer{B}

\fourchoices[answer=B]
{$5\Ums$}
{$10\Ums$}
{$15\Ums$}
{$20\Ums$}

%% answer: B
%% memo:
\memoanswer{
由题知，$v=50\Ums$，$a=6\Umsq$，$x=200\Um$，最小初速度 $v_{0}=\sqrt{v^{2}-2ax}=\sqrt{50^{2}-2\times 6\times 200}\Ums=10\Ums$。故选项B正确。
}

\item
%% number: 14
%% paperName: 2013年广东理综第 14 题 4 分
%% typeId: 1
%% type: 单选题
%% chapter: 曲线运动
%% point: 人造地球卫星
%% method: 无
%% score: 4
%% degree: 650
%% duplicateId: 0
%% body:
如图所示，甲、乙两颗卫星以相同的轨道半径分别绕质量为 $M$ 和 $2M$ 的行星做匀速圆周运动，下列说法正确的是\xzanswer{A}

\onepicture[w=6cm, align=right]{figs/14.png}

\fourchoices[answer=A]
{甲的向心加速度比乙的小}
{甲的运行周期比乙的小}
{甲的角速度比乙大}
{甲的线速度比乙大}

%% answer: A
%% memo:
\memoanswer{
本题原卷及材料中未提供详解，待补充。
}

\item
%% number: 15
%% paperName: 2013年广东理综第 15 题 4 分
%% typeId: 1
%% type: 单选题
%% chapter: 电场
%% point: 带电粒子的偏转
%% method: 无
%% score: 4
%% degree: 650
%% duplicateId: 0
%% body:
喷墨打印机的简化模型如图所示。重力可忽略的墨汁微滴，经带电室带负电后，以速度 $v$ 垂直匀强电场飞入极板间，最终打在纸上，则微滴在极板间电场中\xzanswer{C}

\onepicture[w=7cm, align=right]{figs/15.png}

\fourchoices[answer=C]
{向负极板偏转}
{电势能逐渐增大}
{运动轨迹是抛物线}
{运动轨迹与带电量无关}

%% answer: C
%% memo:
\memoanswer{
本题原卷及材料中未提供详解，待补充。
}

\item
%% number: 16
%% paperName: 2013年广东理综第 16 题 4 分
%% typeId: 1
%% type: 单选题
%% chapter: 交流电
%% point: 变压器
%% method: 无
%% score: 4
%% degree: 650
%% duplicateId: 0
%% body:
如图所示，理想变压器原、副线圈匝数比 $n_{1}:n_{2}=2:1$，V 和 A 均为理想电表，灯泡电阻 $R_{L}=6\UO$，AB 端电压 $u_{1}=12\sqrt{2}\sin100\pi t$（$\UV$）。下列说法正确的是\xzanswer{D}

\onepicture[w=6cm, align=right]{figs/16.png}

\fourchoices[answer=D]
{电流频率为 $100\UHz$}
{V 的读数为 $24\UV$}
{A 的读数为 $0.5\UA$}
{变压器输入功率为 $6\UW$}

%% answer: D
%% memo:
\memoanswer{
本题原卷及材料中未提供详解，待补充。
}

\item
%% number: 17
%% paperName: 2013年广东理综第 17 题 6 分
%% typeId: 2
%% type: 多选题
%% chapter: 原子和原子核
%% point: 核裂变
%% method: 无
%% score: 6
%% degree: 650
%% duplicateId: 0
%% body:
铀核裂变是核电站核能的重要来源，其一种裂变反应 \ce{^{235}_{92}U + ^{1}_{0}n -> ^{144}_{56}Ba + ^{89}_{36}Kr + 3^{1}_{0}n}，下列说法正确的有\xzanswer{AC}

\fourchoices[answer=AC]
{上述裂变反应中伴随着中子放出}
{铀块体积对链式反应的发生无影响}
{铀核的链式反应可人工控制}
{铀核的半衰期会受到环境温度的影响}

%% answer: AC
%% memo:
\memoanswer{
本题原卷及材料中未提供详解，待补充。
}

\item
%% number: 18
%% paperName: 2013年广东理综第 18 题 6 分
%% typeId: 2
%% type: 多选题
%% chapter: 气体性质
%% point: 热力学第一定律
%% method: 无
%% score: 6
%% degree: 650
%% duplicateId: 0
%% body:
如图所示为某同学设计的喷水装置，内部装有 $2\UL$ 水，上部密封 $1\Uatm$ 的空气 $0.5\UL$ 保持阀门关闭，再充入 $1\Uatm$ 的空气 $0.1\UL$，设在所有过程中空气可看作理想气体，且温度不变，下列说法正确的有\xzanswer{AC}

\onepicture[w=5cm, align=right]{figs/18.png}

\fourchoices[answer=AC]
{充气后，密封气体压强增加}
{充气后，密封气体的分子平均动能增加}
{打开阀门后，密封气体对外界做正功}
{打开阀门后，不再充气也能把水喷光}

%% answer: AC
%% memo:
\memoanswer{
本题原卷及材料中未提供详解，待补充。
}

\item
%% number: 19
%% paperName: 2013年广东理综第 19 题 6 分
%% typeId: 2
%% type: 多选题
%% chapter: 功和能
%% point: 机械能守恒定律
%% method: 图象法
%% score: 6
%% degree: 450
%% duplicateId: 0
%% body:
如图所示，游乐场中，从高处 A 到水面 B 处有两条长度相等的光滑轨道。甲、乙两小孩沿着不同轨道同时从 A 处自由滑向 B 处，下列说法正确的有\xzanswer{BD}

\onepicture[w=7cm, align=right]{figs/19.png}

\fourchoices[answer=BD]
{甲的切向加速度始终比乙大}
{甲、乙在同一高度的速度大小相等}
{甲、乙在同一时刻总能到达同一高度}
{甲比乙先到达 B 处}

%% answer: BD
%% memo:
\memoanswer{
本题原卷及材料中未提供详解，待补充。
}

\item
%% number: 20
%% paperName: 2013年广东理综第 20 题 6 分
%% typeId: 2
%% type: 多选题
%% chapter: 力物体的平衡
%% point: 物体系平衡
%% method: 整体与隔离
%% score: 6
%% degree: 650
%% duplicateId: 0
%% body:
如图所示，物体 P 静止于固定的斜面上，P 的上表面水平，现把物体 Q 轻轻地叠放在 P 上，则\xzanswer{BD}

\onepicture[w=7cm, align=right]{figs/20.png}

\fourchoices[answer=BD]
{P 向下滑动}
{P 静止不动}
{P 所受的合外力增大}
{P 与斜面间的静摩擦力增大}

%% answer: BD
%% memo:
\memoanswer{
本题原卷及材料中未提供详解，待补充。
}

\item
%% number: 21
%% paperName: 2013年广东理综第 21 题 6 分
%% typeId: 2
%% type: 多选题
%% chapter: 磁场
%% point: 带电粒子在磁场中的运动
%% method: 无
%% score: 6
%% degree: 450
%% duplicateId: 0
%% body:
如图所示，两个初速度大小相同的同种离子 a 和 b，从 O 点沿垂直磁场方向进入匀强磁场，最后打到屏 P 上。不计重力。下列说法正确的有\xzanswer{AD}

\onepicture[w=4.5cm, align=right]{figs/21.png}

\fourchoices[answer=AD]
{a、b 均带正电}
{a 在磁场中飞行的时间比 b 的短}
{a 在磁场中飞行的路程比 b 的短}
{a 在 P 上的落点与 O 点的距离比 b 的近}

%% answer: AD
%% memo:
\memoanswer{
本题原卷及材料中未提供详解，待补充。
}

\item
%% number: 34
%% paperName: 2013年广东理综第 34 题 9 分
%% typeId: 4
%% type: 实验题
%% chapter: 电路
%% point: 电路实验
%% method: 无
%% score: 9
%% degree: 650
%% duplicateId: 0
%% body:
（1）研究小车匀变速直线运动的实验装置如图（a）所示，其中斜面倾角可调。打点计时器的工作频率为 $50\UHz$。纸带上计数点的间距如图（b）所示，其中每相邻两点之间还有 4 个记录点未画出。
\begin{enumerate}
	\item
	部分实验步骤如下：

	（A）测量完毕，关闭电源，取出纸带。

	（B）接通电源，待打点计时器工作稳定后放开小车。

	（C）将小车停靠在打点计时器附近，小车尾部与纸带相连。

	（D）把打点计时器固定在平板上，让纸带穿过限位孔。

	上述实验步骤的正确顺序是：\tkanswer[4.5cm]{CDBA}（用字母填写）。

	\item
	图（b）中标出的相邻两计数点的时间间隔 $T=$\tkanswer[2cm]{0.1} $\Us$。

	\item
	计数点 5 对应的瞬时速度大小计算式为 $v_{5}=$\tkanswer[3.5cm]{$\frac{s_{4}+s_{5}}{2T}$}。

	\item
	为了充分利用记录数据，减小误差，小车加速度大小的计算式应为 $a=$\tkanswer[4cm]{$\frac{(s_{4}+s_{5}+s_{6})-(s_{1}+s_{2}+s_{3})}{9T^{2}}$}。
\end{enumerate}

\onepicture[w=12cm, align=right]{figs/34.png}

（2）图（a）是测量电阻 $R_{x}$ 的原理图，学生电源输出电压可调，电流表量程选 $0.6\UA$（内阻不计），标有长度刻度的均匀电阻丝 ab 的总长为 $30.0\Ucm$。
\begin{enumerate}
	\item
	根据原理图连接图（b）的实物图。

	\item
	断开 $S_{2}$，合上 $S_{1}$；调节电源输出电压为 $3.0\UV$ 时，单位长度电阻丝的电压 $U=$\tkanswer[1.5cm]{0.1} $\UVcm$。记录此时电流表 $A_{1}$ 的示数。

	\item
	保持 $S_{1}$ 闭合，合上 $S_{2}$；滑动 c 点改变 ac 的长度 $L$，同时调节电源输出电压，使电流表 $A_{1}$ 的示数与步骤（2）记录的值相同，记录长度 $L$ 和 $A_{2}$ 的示数 $I$。测量 6 组 $L$ 和 $I$ 值，测量数据已在图（c）中标出。写出 $R_{x}$ 与 $L$、$I$、$U$ 的关系式 $R_{x}=$\tkanswer[2.5cm]{$\frac{LU}{I}$}；根据图（c）用作图法算出 $R_{x}=$\tkanswer[1.5cm]{6.0} $\UO$。
\end{enumerate}

\threepicture[align=right, showlabel=false, labelA=2013广东34a, labelB=2013广东34b, labelC=2013广东34c, hA=4.5cm, hB=6cm, hC=6cm]
{figs/34a.png}
{figs/34b.png}
{figs/34c.png}

%% answer:
\jdanswer{
\begin{enumerate}
	\item
	CDBA

	\item
	$0.1\Us$

	\item
	$\frac{s_{4}+s_{5}}{2T}$

	\item
	$\frac{(s_{4}+s_{5}+s_{6})-(s_{1}+s_{2}+s_{3})}{9T^{2}}$

	\item
	见解析（实物连线如图所示）

	\onepicture[w=8cm]{figs/34b_an.png}

	\item
	0.1

	\item
	$\frac{LU}{I}$，6.0（描点连线图线如图所示）

	\onepicture[h=7.5cm]{figs/34c_an.png}

\end{enumerate}
}
%% memo:
\memoanswer{
本题原卷及材料中未提供详解，待补充。
}

\item
%% number: 35
%% paperName: 2013年广东理综第 35 题 18 分
%% typeId: 6
%% type: 计算题
%% chapter: 运动和力
%% point: 动量守恒定律
%% method: 无
%% score: 18
%% degree: 650
%% duplicateId: 0
%% body:
如图所示，两块相同平板 $P_{1}$、$P_{2}$ 置于光滑水平面上，质量均为 $m$。$P_{2}$ 的右端固定一轻质弹簧，左端 A 与弹簧的自由端 B 相距 $L$。物体 P 置于 $P_{1}$ 的最右端，质量为 $2m$ 且可看作质点。$P_{1}$ 与 P 以共同速度 $v_{0}$ 向右运动，与静止的 $P_{2}$ 发生碰撞，碰撞时间极短，碰撞后 $P_{1}$ 与 $P_{2}$ 粘连在一起。P 压缩弹簧后被弹回并停在 A 点（弹簧始终在弹性限度内）。P 与 $P_{2}$ 之间的动摩擦因数为 $\mu$。求：
\begin{enumerate}
	\item
	$P_{1}$、$P_{2}$ 刚碰完时的共同速度 $v_{1}$ 和 P 的最终速度 $v_{2}$；
	\item
	此过程中弹簧的最大压缩量 $x$ 和相应的弹性势能 $E_{p}$。
\end{enumerate}

\onepicture[w=9cm, align=right]{figs/35.png}

%% answer:
\jdanswer{
\begin{enumerate}
	\item
	$v_{1}=\frac{1}{2}v_{0}$，$v_{2}=\frac{3}{4}v_{0}$

	\item
	$x=\frac{v_{0}^{2}}{32\mu}-L$，$E_{p}=\frac{mv_{0}^{2}}{16}$

\end{enumerate}
}
%% memo:
\memoanswer{
（1）$P_{1}$ 和 $P_{2}$ 碰撞动量守恒：
$mv_{0}=(m+m)v_{1}$ ①

得出：$v_{1}=\frac{1}{2}v_{0}$。

P 在 $P_{2}$ 上滑行过程 $P_{1}$、$P_{2}$、P 系统动量守恒：
$2mv_{0}+2mv_{1}=4mv_{2}$ ②

得出：$v_{2}=\frac{3}{4}v_{0}$。

（2）$P_{1}$、$P_{2}$、P 第一次等速弹簧最大压缩量最大，由能量守恒得
$\mu\cdot 2mg(L+x)+E_{p}=\frac{1}{2}(2m)v_{0}^{2}+\frac{1}{2}(2m)v_{1}^{2}-\frac{1}{2}(4m)v_{2}^{2}$ ③

P 刚进入 $P_{2}$ 到 $P_{1}$、$P_{2}$、P 第二次等速时有能量守恒得：
$\mu\cdot 2mg(2L+2x)=\frac{1}{2}(2m)v_{0}^{2}+\frac{1}{2}(2m)v_{1}^{2}-\frac{1}{2}(4m)v_{2}^{2}$ ④

由③④得：$x=\frac{v_{0}^{2}}{32\mu}-L$，$E_{p}=\frac{mv_{0}^{2}}{16}$。
}

\item
%% number: 36
%% paperName: 2013年广东理综第 36 题 18 分
%% typeId: 6
%% type: 计算题
%% chapter: 电磁感应
%% point: 电磁感应电路分析
%% method: 无
%% score: 18
%% degree: 450
%% duplicateId: 0
%% body:
如图（a）所示，在垂直于匀强磁场 B 的平面内，半径为 $r$ 的金属圆盘绕过圆心 O 的轴承转动，圆心 O 和边缘 K 通过电刷与一个电路连接。电路中的 P 是加上一定正向电压才能导通的电子元件。流过电流表的电流 $I$ 与圆盘角速度 $\omega$ 的关系如图（b）所示，其中 ab 段和 bc 段均为直线，且 ab 段过坐标原点。$\omega>0$ 代表圆盘逆时针转动。已知：$R=3.0\UO$，$B=1.0\UT$，$r=0.2\Um$。忽略圆盘，电流表和导线的电阻。
\begin{enumerate}
	\item
	根据图（b）写出 ab、bc 段对应的 $I$ 与 $\omega$ 的关系式；
	\item
	求出图（b）中 b、c 两点对应的 P 两端的电压 $U_{b}$、$U_{c}$；
	\item
	分别求出 ab、bc 段流过 P 的电流 $I_{P}$ 与其两端电压 $U_{P}$ 的关系式。
\end{enumerate}

\onepicture[w=12cm, align=right]{figs/36.png}

%% answer:
\jdanswer{
\begin{enumerate}
	\item
	当 $-45\leq\omega\leq 15$ 时，$I=\frac{\omega}{150}\UA$；当 $15\leq\omega\leq 45$ 时，$I=\left(\frac{\omega}{100}-0.05\right)\UA$

	\item
	$U_{b}=0.3\UV$，$U_{c}=0.9\UV$

	\item
	当 $-45\leq\omega\leq 15$ 时，$U_{p}=I_{p}R_{p}=0\UV$；当 $15\leq\omega\leq 45$ 时，$U_{p}=I_{p}R_{p}=9\left(\frac{\omega}{300}-0.05\right)\UV$

\end{enumerate}
}
%% memo:
\memoanswer{
（1）由图像得出三点坐标 $O(0,0)$、$b(15,0.1)$、$c(45,0.4)$，由直线的两点式得 $I$ 与 $\omega$ 的关系式：当 $-45\leqslant\omega\leqslant 15$ 时，$I=\frac{\omega}{150}\UA$；当 $15\leqslant\omega\leqslant 45$ 时，$I=\left(\frac{\omega}{100}-0.05\right)\UA$。

（2）圆盘切割产生的电动势为：$E=Br\frac{\omega r+0}{2}=\frac{1}{2}B\omega r^{2}=0.02\omega$。当 $\omega=15$ 时 $E=0.3\UV$，当 $\omega=45$ 时 $E=0.9\UV$。电源忽略内阻，故 $U_{p}=E$，可得：$U_{b}=0.3\UV$，$U_{c}=0.9\UV$。

（3）由并联电路知识有：$I=I_{p}+I_{R}$ ①，$I_{R}=\frac{E}{R}$ ②。由①②得 $I_{p}=I-\frac{E}{R}=I-\frac{0.02\omega}{3}=I-\frac{\omega}{150}$，即 $I_{p}=\begin{cases}0, & -45\leqslant\omega\leqslant 15\\ \left(\frac{\omega}{300}-0.05\right)\UA, & 15<\omega\leqslant 45\end{cases}$。由 $\omega=45$ 时 $E=0.9\UV$，C 点导通，$I=I_{p}+I_{R}=\frac{E}{R_{p}}+\frac{E}{R}$，且 $R=3\UO$，$I=0.4\UA$，得 $R_{p}=9\UO$。当 $-45\leqslant\omega\leqslant 15$ 时，$U_{p}=I_{p}R_{p}=0\UV$；当 $15\leqslant\omega\leqslant 45$ 时，$U_{p}=I_{p}R_{p}=9\left(\frac{\omega}{300}-0.05\right)\UV$。
}

\end{enumerate}

\end{document}
